\section*{Упражнения}
\addcontentsline{toc}{section}{Упражнения}

\begin{exercise}[2-1]
Пусть $F$~--- поле характеристики $\ne2$.
\begin{enumerate}[label=(\alph*)]
\item Пусть $E$~--- квадратичное расширение $F$. Покажите, что
\[
S(E)=\{a\in F^\times\mid a\ \text{является квадратом в}\ E\}
\]
~--- подгруппа $F^\times$, содержащая $F^{\times2}$.
\item Пусть $E$ и $E'$~--- квадратичные расширения $F$. Покажите, что $F$-изоморфизм $\varphi\colon E\to E'$ существует
тогда и только тогда, когда $S(E)=S(E')$.
\item Покажите, что существует бесконечная последовательность полей $E_1,E_2,\dots$, являющихся квадратичными
расширениями $\Q$, в которой $E_i$ не изоморфно $E_j$ при $i\ne j$.
\item Пусть $p$~--- нечётное простое число. Покажите, что с точностью до изоморфизма существует ровно одно поле из
$p^2$ элементов.
\end{enumerate}
\end{exercise}

\noindent\emph{Решение.} (a) Если $a=x^2$, $b=y^2$ с $x,y\in E$, то $ab=(xy)^2$ и $a^{-1}=(x^{-1})^2$; кроме того,
$1=1^2$. Значит, $S(E)$~--- подгруппа. Если $c\in F^\times$, то $c^2$~--- квадрат элемента $c\in E$, т.\,е.\
$F^{\times2}\subset S(E)$.

Для (b)--(d) сначала опишем квадратичные расширения и группу $S(E)$ явно.

\emph{Лемма.} Пусть $[E:F]=2$. Тогда $E=F[\delta]$, где $\delta^2=d\in F^\times\smallsetminus F^{\times2}$, и
\[
S(E)=F^{\times2}\cup dF^{\times2},\qquad (S(E):F^{\times2})=2 .
\]

Действительно, возьмём $\alpha\in E\smallsetminus F$. Тогда $F\subsetneq F[\alpha]\subset E$, и по \rf{1.20}
$[F[\alpha]:F]$ делит $2$ и больше $1$; значит, $[F[\alpha]:F]=2$, $[E:F[\alpha]]=1$ и $E=F[\alpha]$ (собственного
промежуточного поля нет). Пусть $X^2+bX+c$~--- минимальный многочлен $\alpha$. Положим $\delta=2\alpha+b$ (<<выделение
полного квадрата>>):
\[
\delta^2=4\alpha^2+4b\alpha+b^2=4(\alpha^2+b\alpha)+b^2=b^2-4c=:d\in F .
\]
Так как $2$ обратимо в $F$, то $\delta\notin F$ и $F[\delta]=F[\alpha]=E$; в частности, $d\ne0$. Элемент $d$ не квадрат в
$F$: если $d=e^2$, $e\in F$, то $\delta$, как корень $X^2-d=(X-e)(X+e)$, равнялся бы $\pm e\in F$. Далее, $1,\delta$~---
базис $E$ над $F$. Если $a\in F^\times$ и $a=(x+y\delta)^2=x^2+dy^2+2xy\delta$ с $x,y\in F$, то $2xy=0$, т.\,е.\ $y=0$ и
$a=x^2\in F^{\times2}$, либо $x=0$ и $a=dy^2\in dF^{\times2}$. Обратно, $x^2=(x)^2$ и $dy^2=(y\delta)^2$ лежат в $S(E)$.
Смежные классы $F^{\times2}$ и $dF^{\times2}$ различны, так как $d\notin F^{\times2}$; поэтому индекс равен $2$.

(b) Пусть $\varphi\colon E\to E'$~--- $F$-изоморфизм. Если $a\in S(E)$, $a=x^2$, то $a=\varphi(a)=\varphi(x)^2$, т.\,е.\
$a\in S(E')$. Итак, $S(E)\subset S(E')$, и, применяя то же к $\varphi^{-1}$, получаем $S(E)=S(E')$. (Это направление
верно для любых расширений $E,E'$, не только квадратичных.)

Обратно, пусть $S(E)=S(E')$. По лемме $E=F[\delta]$, $\delta^2=d$. Так как $d\in S(E)=S(E')$, найдётся $\delta'\in E'$ с
$\delta'^2=d$. Многочлен $X^2-d$ неприводим над $F$ (у него нет корней в $F$) и является минимальным многочленом
$\delta$; поэтому по \rf{2.1}(b) существует $F$-гомоморфизм $\varphi\colon F[\delta]\to E'$, $\delta\mapsto\delta'$.
Он инъективен (как всякий гомоморфизм полей) и $F$-линеен, а $\dim_FE=\dim_FE'=2$; инъективное линейное отображение
пространств одинаковой конечной размерности биективно. Значит, $\varphi$~--- $F$-изоморфизм.

\par\noindent\begin{minipage}{\linewidth}
\begin{center}
\begin{tikzcd}[column sep=large,row sep=small]
& E=F[\delta]\arrow[dd,cnew,"\varphi"',"\delta\,\mapsto\,\delta'"]\\
F\arrow[ur,cold,hook]\arrow[dr,cold,hook] &\\
& E'\ni\delta',\ \delta'^2=d
\end{tikzcd}
\end{center}
\figcap{Из $d\in S(E')$ получаем корень $\delta'\in E'$ многочлена $X^2-d$, а по нему~--- $F$-изоморфизм $\varphi$.}
\end{minipage}\par\medskip

Итак, \emph{квадратичные расширения $F$ с точностью до $F$-изоморфизма взаимно однозначно соответствуют подгруппам
порядка $2$ группы $F^\times/F^{\times2}$}, т.\,е.\ её неединичным элементам (каждый элемент этой группы имеет порядок
$\le2$): расширению $E=F[\sqrt d]$ отвечает подгруппа $S(E)/F^{\times2}=\{1,\bar d\}$. Соответствие инъективно по (b) и
сюръективно: для $d\notin F^{\times2}$ поле $E=F[X]/(X^2-d)$ квадратично, класс $x$ переменной $X$ удовлетворяет
$x^2=d$, т.\,е.\ $d\in S(E)\smallsetminus F^{\times2}$, и по лемме $S(E)/F^{\times2}=\{1,\bar d\}$.
Например, $\R^\times/\R^{\times2}=\{\pm1\}$, и у $\R$ одно квадратичное расширение, $\CC$; у $\CC$ их нет. Это случай
$n=2$ теории Куммера (\rf{5.30}; группа $S(E)$~--- это группа $B(E)$ из \rf{5.30}): конечные подгруппы
$F^\times/F^{\times2}$ большего порядка описывают мультиквадратичные расширения $F[\sqrt{a_1},\dots,\sqrt{a_r}]$;
например, полю $\Q[\sqrt2,\sqrt3]$ отвечает подгруппа $\{1,2,3,6\}$ (по модулю $\Q^{\times2}$, \rf{5.31a}).

(c) Пусть $p$ пробегает простые числа, $E_p=\Q[\sqrt p]$ ($X^2-p$ неприводим по Эйзенштейну, \rf{1.16a}(a)). По лемме
$S(E_p)=\Q^{\times2}\cup p\Q^{\times2}$. Пусть $q\ne p$~--- другое простое. Тогда $q\notin S(E_p)$: в разложении квадрата
рационального числа на простые множители все показатели чётны, а у $q$ и у $q/p$ показатель при $q$ равен $1$. Значит,
$S(E_p)\ne S(E_q)$, и по (b) поля $E_p$ и $E_q$ не $\Q$-изоморфны. Но любой изоморфизм полей характеристики $0$ тождествен
на $\Q$ (он переводит $1$ в $1$, а значит, $n\cdot1$ в $n\cdot1$ и $m/n$ в $m/n$), т.\,е.\ является $\Q$-изоморфизмом.
Поэтому поля $E_2,E_3,E_5,E_7,\dots$ попарно неизоморфны.

То же можно проверить непосредственно: если $\sqrt p=a+b\sqrt q$, $a,b\in\Q$, то $p=a^2+qb^2+2ab\sqrt q$, откуда $ab=0$
(иначе $\sqrt q\in\Q$). При $b=0$ получаем $\sqrt p=a\in\Q$, а при $a=0$, $b=m/n$,~--- равенство $pn^2=qm^2$, в котором
показатель $p$ слева нечётен, а справа чётен. Значит, $\sqrt p\notin E_q$, хотя $\Q$-изоморфизм $E_p\to E_q$ перевёл бы
$\sqrt p$ в корень $X^2-p$ в $E_q$. Вообще, $\Q^\times/\Q^{\times2}$~--- векторное пространство над $\F_2$ с базисом из
классов $-1,2,3,5,7,\dots$ (разложение на простые множители), так что квадратичных расширений $\Q$ бесконечно много.

(d) Пусть $F$~--- поле из $p^2$ элементов. Его характеристика не равна $0$ (иначе $F\supset\Q$ бесконечно), т.\,е.\ равна
простому $\ell$, и $F$~--- векторное пространство над $\F_\ell$ некоторой размерности $n$; тогда $p^2=|F|=\ell^n$, так что
$\ell=p$, $n=2$: $F$~--- квадратичное расширение $\F_p$.

\emph{Существование.} Квадраты образуют подгруппу индекса $2$ в $\F_p^\times$ (\rf{1.18b}); для неквадрата $d$ поле
$\F_p[X]/(X^2-d)$ содержит $p^2$ элементов.

\emph{Единственность.} Для любого квадратичного расширения $F/\F_p$ по лемме $(S(F):\F_p^{\times2})=2=(\F_p^\times:\F_p^{\times2})$,
поэтому $S(F)=\F_p^\times$: \emph{в любом поле из $p^2$ элементов каждый элемент $\F_p$~--- квадрат}. Значит, все $S(F)$
одинаковы, и по (b) все такие поля $\F_p$-изоморфны (а любой изоморфизм тождествен на простом подполе $\F_p$). Заметим,
что группа $S(F)$ всегда строго больше $\F_p^{\times2}$: по лемме она содержит неквадрат $d$.

Пример: в $\F_7$ квадраты~--- это $\{1,2,4\}$, неквадраты~--- $\{3,5,6\}$. В поле $\F_7[\sqrt3]$ имеем
$(2\sqrt3)^2=12=5$ и $(3\sqrt3)^2=27=6$, так что $\F_7[\sqrt5]=\F_7[\sqrt6]=\F_7[\sqrt3]$: произведение двух неквадратов~---
квадрат (\rf{1.18b}), и $\sqrt5=\sqrt{5/3}\cdot\sqrt3$ с $5/3=4=2^2$.

Другое объяснение того, что $S(F)=\F_p^\times$: группа $F^\times$ циклична порядка $p^2-1$ (упражнение 1-3), и если
$g$~--- её образующая, то единственная подгруппа порядка $p-1$, т.\,е.\ $\F_p^\times$, состоит из степеней $g^{(p+1)k}$, а
$p+1$ чётно. Для $p=2$ утверждение (d) тоже верно (поле из $4$ элементов~--- это $\F_2[X]/(X^2+X+1)$, так как
$X^2+X+1$~--- единственный неприводимый квадратный многочлен над $\F_2$), но рассуждение с $S(F)$ не работает: над $\F_2$
каждый элемент~--- квадрат, и квадратичные расширения в характеристике $2$ устроены иначе (они получаются из многочленов
$X^2+X+a$, \rf{2.22b}, или $X^2-a$, \rf{2.18}). Общий результат о полях из $p^m$ элементов~--- в упражнении 2-4 и в
\rf{4.20}.

\begin{exercise}[2-2]
\begin{enumerate}[label=(\alph*)]
\item Пусть $F$~--- поле характеристики $p$. Покажите, что если $X^p-X-a$ приводим в $F[X]$, то он раскладывается в
$F[X]$ на различные множители.
\item Покажите, что для каждого простого $p$ многочлен $X^p-X-1$ неприводим в $\Q[X]$.
\end{enumerate}
\end{exercise}

\noindent\emph{Решение.} (a) Это \rf{2.22b}(b): если $X^p-X-a$ приводим, то $a=\beta^p-\beta$ для некоторого $\beta\in F$,
и $X^p-X-a=\prod_{i\in\F_p}(X-\beta-i)$~--- произведение $p$ различных линейных (и, значит, неприводимых) множителей.
Напомним идею: если $g$~--- унитарный делитель степени $0<r<p$, то $g=\prod_{i\in I}(X-\alpha-i)$ для корня $\alpha$, и
коэффициент $-r\alpha-\sum_{i\in I}i\in F$ при $X^{r-1}$ даёт $\alpha\in F$.

(b) Пусть $f=X^p-X-1\in\Z[X]$. Его приведение $\bar f=X^p-X-1\in\F_p[X]$ неприводимо по (a): для $\beta\in\F_p$ имеем
$\beta^p-\beta=0\ne1$ (малая теорема Ферма), т.\,е.\ у $\bar f$ нет корней в $\F_p$, а приводимый многочлен вида
$X^p-X-a$ имел бы их по (a). Старший коэффициент $f$ равен $1$ и не делится на $p$; по \rf{1.18} многочлен $f$ неприводим
в $\Q[X]$.

Сдвиг $X\mapsto X+1$, превращающий $\Phi_p$ в многочлен Эйзенштейна (\rf{1.16a}(c)), здесь не помогает:
$f(X+1)=X^p+\sum_{k=2}^{p-1}\binom pkX^k+(p-1)X-1\equiv f(X)\pmod p$, и вообще $f(X+c)\equiv f(X)\pmod p$ при $c\in\Z$
(так как $c^p\equiv c$); свободный член $f(c)\equiv-1$ не делится на $p$, и ни один целый сдвиг не даёт многочлена
Эйзенштейна относительно $p$. Это то же наблюдение, что в \rf{2.9}:
сдвиг на $1$ переставляет корни $\bar f$.

\begin{exercise}[2-3]
Постройте поле разложения многочлена $X^5-2$ над $\Q$. Какова его степень над $\Q$?
\end{exercise}

\noindent\emph{Решение.} Ответ: $E=\Q[\alpha,\zeta]$, где $\alpha=\sqrt[5]2\in\R$ и $\zeta=e^{2\pi i/5}$;
$[E:\Q]=20$.

Корни $X^5-2$ в $\CC$~--- это $\alpha\zeta^k$, $k=0,\dots,4$; они попарно различны, и
\[
X^5-2=(X-\alpha)(X-\alpha\zeta)(X-\alpha\zeta^2)(X-\alpha\zeta^3)(X-\alpha\zeta^4).
\]
Поле разложения внутри $\CC$~--- это $\Q[\alpha,\alpha\zeta,\dots,\alpha\zeta^4]$. Оно содержит $\alpha$ и
$\zeta=\alpha\zeta/\alpha$, а $\Q[\alpha,\zeta]$ содержит все корни; значит, поле разложения равно $\Q[\alpha,\zeta]$:
оно обязано содержать $\alpha$ и $\zeta$, и этого достаточно. (Без $\CC$ его можно построить, как в \rf{2.7}: поле корня
$\Q[\alpha]\cong\Q[X]/(X^5-2)$, затем поле корня неприводимого множителя многочлена $(X^5-2)/(X-\alpha)$ над $\Q[\alpha]$; ниже
мы увидим, что этот множитель~--- весь многочлен четвёртой степени. По \rf{2.13} результат $\Q$-изоморфен
$\Q[\alpha,\zeta]$.)

\emph{Степень.} Многочлен $X^5-2$ неприводим (Эйзенштейн при $p=2$, \rf{1.16a}(a)), поэтому $[\Q[\alpha]:\Q]=5$. Число
$\zeta$~--- корень многочлена $\Phi_5=X^4+X^3+X^2+X+1$, неприводимого над $\Q$ (\rf{1.16a}(c)), поэтому
$[\Q[\zeta]:\Q]=4$. По \rf{1.20} оба числа $5$ и $4$ делят $[E:\Q]$, т.\,е.\ $20\mid[E:\Q]$. С другой стороны,
$E=\Q[\zeta][\alpha]$, и $\alpha$~--- корень многочлена $X^5-2$ степени $5$ над $\Q[\zeta]$, так что $[E:\Q[\zeta]]\le5$ и
$[E:\Q]\le5\cdot4=20$. Итак, $[E:\Q]=20$.

\par\noindent\begin{minipage}{\linewidth}
\begin{center}
\begin{tikzpicture}[every node/.style={font=\small}]
\node (Q) at (0,0) {$\Q$};
\node[cold] (A) at (-1.6,1.3) {$\Q[\alpha]$};
\node[cold] (Z) at (1.6,1.1) {$\Q[\zeta]$};
\node[cnew] (E) at (0,2.6) {$E=\Q[\alpha,\zeta]$};
\draw (Q) -- node[below left=-2pt] {$5$} (A);
\draw (Q) -- node[below right=-2pt] {$4$} (Z);
\draw (A) -- node[above left=-2pt] {$4$} (E);
\draw (Z) -- node[above right=-2pt] {$5$} (E);
\end{tikzpicture}
\end{center}
\figcap{Степени в башнях: $20=5\cdot4=4\cdot5$.}
\end{minipage}\par\medskip

Попутно получаем: $[E:\Q[\zeta]]=5$, т.\,е.\ $X^5-2$ остаётся неприводимым над $\Q[\zeta]$, и $[E:\Q[\alpha]]=4$. Над
$\Q[\alpha]$
\[
X^5-2=(X-\alpha)\bigl(X^4+\alpha X^3+\alpha^2X^2+\alpha^3X+\alpha^4\bigr),\qquad
X^4+\alpha X^3+\alpha^2X^2+\alpha^3X+\alpha^4=\alpha^4\,\Phi_5(X/\alpha),
\]
и второй множитель неприводим над $\Q[\alpha]$: его корень $\alpha\zeta$ порождает $E$ над $\Q[\alpha]$, и степень этого
корня над $\Q[\alpha]$ равна $[E:\Q[\alpha]]=4$. Аналогично $\Phi_5$ неприводим над $\Q[\alpha]$.

\begin{exercise}[2-4]
Найдите поле разложения многочлена $X^{p^m}-1\in\F_p[X]$. Какова его степень над $\F_p$?
\end{exercise}

\noindent\emph{Решение.} Возведение в степень $p^m$ в кольце $\F_p[X]$ характеристики $p$ аддитивно, поэтому
\[
X^{p^m}-1=(X-1)^{p^m}.
\]
Многочлен уже раскладывается над $\F_p$: его поле разложения~--- само $\F_p$, степень равна $1$; единственный корень $1$
имеет кратность $p^m$. Это частный случай \rf{2.10c}: поле разложения $X^N-1$ над $\F_p$ имеет степень, равную порядку
$p$ по модулю наибольшего делителя $N$, взаимно простого с~$p$, а здесь этот делитель равен $1$; в характеристике $p$ нет корней из единицы степени
$p$, кроме $1$ (\rf{2.10}(c)).

Интереснее похожий многочлен $X^{p^m}-X$. Положим $q=p^m$, и пусть $E$~--- поле разложения $X^q-X$ над $\F_p$.
\begin{itemize}
\item $X^q-X$ сепарабелен: его производная $qX^{q-1}-1=-1$. Поэтому у него ровно $q$ различных корней в $E$.
\item Множество $R$ этих корней~--- подполе $E$: отображение $x\mapsto x^q$~--- гомоморфизм колец (композиция $m$
эндоморфизмов Фробениуса), так что из $x^q=x$, $y^q=y$ следует $(x\pm y)^q=x\pm y$, $(xy)^q=xy$ и
$(x^{-1})^q=x^{-1}$ при $x\ne0$; кроме того, $0,1\in R$.
\item Поле $E$ порождено корнями, т.\,е.\ $E=\F_p[R]=R$. Значит, $|E|=q=p^m$ и $[E:\F_p]=m$.
\end{itemize}
Итак, \emph{поле разложения $X^{p^m}-X$ над $\F_p$~--- это поле из $p^m$ элементов, его степень равна $m$.}
Обратно, если $K$~--- любое поле из $q$ элементов, то $|K^\times|=q-1$, и по теореме Лагранжа $x^{q-1}=1$ для всех
$x\in K^\times$, т.\,е.\ $x^q=x$ для всех $x\in K$. Поэтому $X^q-X=\prod_{x\in K}(X-x)$ раскладывается в $K$, и $K$
порождено его корнями~--- это поле разложения $X^q-X$. По \rf{2.13} \emph{любые два поля из $p^m$ элементов изоморфны}
(ср.\ \rf{4.20}); при $m=2$ это ещё раз доказывает 2-1(d).

Два дополнения. (1) При $d\mid m$ многочлен $X^{p^d}-X$ делит $X^{p^m}-X$: число $p^d-1$ делит $p^m-1=(p^d)^{m/d}-1$, а
$X^k-1$ делит $X^{kl}-1$; остаётся умножить на $X$. В частности, $X^p-X=\prod_{k\in\F_p}(X-k)$ делит $X^{p^m}-X$, и
вообще $p^d$ корней $X^{p^d}-X$ образуют в $E$ подполе из $p^d$ элементов (по тому же рассуждению, что для $R$; ср.\
\rf{4.21}). (2) Поле $K$ из $p^m$ элементов имеет вид $\F_p[X]/(f)$ с неприводимым $f$ степени $m$: группа $K^\times$
циклична (упражнение 1-3(e)), её образующая $\gamma$ порождает $K$ над $\F_p$, минимальный многочлен $f$ элемента
$\gamma$ имеет степень $[K:\F_p]=m$, и $K=\F_p[\gamma]\cong\F_p[X]/(f)$ (\rf{1.27c}). Так как поле из $p^m$ элементов
существует, над $\F_p$ есть неприводимые многочлены любой степени $m\ge1$. Обратно, каждый унитарный неприводимый $f$
степени $m$ делит $X^{p^m}-X$: в поле $\F_p[X]/(f)$ из $p^m$ элементов класс $x$ удовлетворяет $x^{p^m}=x$, а $f$~---
минимальный многочлен $x$. Полное разложение $X^{p^m}-X$ на неприводимые множители~--- в \rf{4.22b}.

\begin{exercise}[2-5]
Пусть $f\in F[X]$, где $F$~--- поле характеристики $0$, и $d(X)=\gcd(f,f')$. Покажите, что многочлен
$g(X)=f(X)d(X)^{-1}$ имеет те же корни, что и $f(X)$, и все они простые корни $g(X)$.
\end{exercise}

\noindent\emph{Решение.} Будем считать $f$ непостоянным. Пусть в поле разложения $E$
\[
f=a\prod_{i=1}^r(X-\alpha_i)^{m_i},\qquad \alpha_i\ \text{попарно различны}.
\]
\emph{Утверждение: $d=\prod_{i=1}^r(X-\alpha_i)^{m_i-1}$.} Запишем $f=(X-\alpha_i)^{m_i}h_i$, $h_i(\alpha_i)\ne0$. По
\rf{2.18a}
\[
f'=(X-\alpha_i)^{m_i-1}\bigl(m_ih_i+(X-\alpha_i)h_i'\bigr),
\]
и второй множитель в точке $\alpha_i$ равен $m_ih_i(\alpha_i)\ne0$, так как $m_i\cdot1_F\ne0$ в характеристике $0$.
Значит, $(X-\alpha_i)^{m_i-1}$ делит $f'$, а $(X-\alpha_i)^{m_i}$ не делит. Унитарный общий делитель $f$ и $f'$ в $E[X]$
делит $f$ и потому имеет вид $\prod(X-\alpha_i)^{k_i}$ с $k_i\le m_i$, а делимость $f'$ даёт $k_i\le m_i-1$; значит, $\gcd(f,f')$ делит $\prod(X-\alpha_i)^{m_i-1}$. Обратно,
многочлены $(X-\alpha_i)^{m_i-1}$ при разных $i$ попарно взаимно просты, поэтому (однозначность разложения в $E[X]$) их
произведение делит $f'$; оно делит и $f$, т.\,е.\ это общий делитель. Поэтому
$\gcd(f,f')$ в $E[X]$ равен $\prod(X-\alpha_i)^{m_i-1}$, и по \rf{2.17} это и есть $d$.

Следовательно,
\[
g=f/d=a\prod_{i=1}^r(X-\alpha_i)\in F[X]
\]
(частное лежит в $F[X]$, так как $d\mid f$ в $F[X]$): корни $g$~--- это в точности $\alpha_1,\dots,\alpha_r$, и все они
простые. Замечательно, что $g$ вычисляется алгоритмом Евклида в $F[X]$, без нахождения корней. Например, для
$f=X^5-X^4-2X^3+2X^2+X-1=(X-1)^3(X+1)^2$ над $\Q$ получаем $d=(X-1)^2(X+1)=X^3-X^2-X+1$ и $g=X^2-1$.

Короче то же можно получить сразу для всех корней. По правилу Лейбница (\rf{2.18a})
\[
f'=a\sum_{i=1}^rm_i(X-\alpha_i)^{m_i-1}\prod_{j\ne i}(X-\alpha_j)^{m_j}=D\cdot S,
\]
где
\[
D=\prod_{i=1}^r(X-\alpha_i)^{m_i-1},\qquad S=a\sum_{i=1}^rm_i\prod_{j\ne i}(X-\alpha_j)
\]
(формально $f'/f=\sum_im_i/(X-\alpha_i)$), а $f=D\cdot a\prod_i(X-\alpha_i)$. Так как $S(\alpha_i)=am_i\prod_{j\ne
i}(\alpha_i-\alpha_j)\ne0$ (снова нужно $m_i\cdot1_F\ne0$), многочлены $S$ и $\prod_i(X-\alpha_i)$ не имеют общих корней,
т.\,е.\ взаимно просты, и в $E[X]$ имеем $\gcd(f,f')=D\cdot\gcd\bigl(\prod_i(X-\alpha_i),S\bigr)=D$.

В характеристике $p$ утверждение верно, если ни одна кратность $m_i$ не делится на $p$ (рассуждение то же), но в общем
случае неверно. Над $\F_p$ для $f=X^{p+1}-X^p=X^p(X-1)$ имеем $f'=(p+1)X^p-pX^{p-1}=X^p$, $d=X^p$ и $g=X-1$: корень $0$
потерян. А для $f=X^p-T$ над $\F_p(T)$ имеем $f'=0$, $d=f$ и $g=1$.

\begin{exercise}[2-6]
Пусть $f(X)$~--- неприводимый многочлен из $F[X]$, где $F$ имеет характеристику $p$. Покажите, что $f$ можно записать в
виде $f(X)=g(X^{p^e})$, где $g(X)$ неприводим и сепарабелен. Выведите, что каждый корень $f(X)$ имеет одну и ту же
кратность $p^e$ в любом поле разложения.
\end{exercise}

\noindent\emph{Решение.} Пусть $e\ge0$~--- наибольшее целое число, для которого $f$ является многочленом от
$X^{p^e}$ (такое есть: при $e=0$ это верно, а $p^e\le\deg f$, если $f\in F[X^{p^e}]$ непостоянен). Запишем
$f(X)=g(X^{p^e})$, $g\in F[X]$.

\emph{$g$ неприводим:} разложение $g=uv$ на множители положительной степени дало бы разложение
$f=u(X^{p^e})\,v(X^{p^e})$ на множители положительной степени.

\emph{$g$ сепарабелен:} $g$ не является многочленом от $X^p$, иначе $f$ был бы многочленом от $X^{p^{e+1}}$ вопреки
выбору $e$. По \rf{2.20} неприводимый $g$ сепарабелен.

\emph{Кратности.} Пусть $g=c\prod_{i=1}^s(X-b_i)$ с попарно различными $b_i$ в некотором расширении $E_1$ (\rf{2.7}), и
пусть в расширении $E_2\supset E_1$ имеются $\beta_i$ с $\beta_i^{p^e}=b_i$ (корни многочленов $X^{p^e}-b_i$). Так как
$x\mapsto x^{p^e}$~--- гомоморфизм,
\[
f(X)=g(X^{p^e})=c\prod_{i=1}^s\bigl(X^{p^e}-\beta_i^{p^e}\bigr)=c\prod_{i=1}^s(X-\beta_i)^{p^e},
\]
а $\beta_i$ попарно различны, так как различны их $p^e$-е степени $b_i$ (и Фробениус инъективен). Значит, $f$ имеет
$s=\deg g$ различных корней, каждый кратности $p^e$, и $\deg f=p^e\deg g$. По независимости набора кратностей от поля
разложения (см.\ подраздел <<Определение кратности и её независимость от поля>>) это верно в любом поле, где $f$ раскладывается.

Число $e$ определено однозначно (это показатель кратности корней), и $g$ тоже: если $f=g_1(X^{p^{e}})$, то $g_1=g$.
Если $f$ сепарабелен, то $e=0$, $g=f$; над совершенным полем так для всех неприводимых $f$ (\rf{2.24}). Примеры:
$X^p-T$ над $\F_p(T)$ ($e=1$, $g=X-T$, один корень кратности $p$); $X^{p^2}-T$ над $\F_p(T)$ ($e=2$, $g=X-T$, один
корень кратности $p^2$; неприводимость~--- по \rf{2.3a} или по Эйзенштейну, \rf{1.16d}(a)); $X^{2p}+tX^p+t$ над $k(t)$
($e=1$, $g=X^2+tX+t$, два корня кратности $p$, пример \rf{2.20a}). Если $\beta$~--- корень $f$, то $\beta^{p^e}$~---
корень сепарабельного $g$; это начало теории сепарабельной степени и степени несепарабельности (\rf{5.46a}).
